\section{Discussion and Conclusion}
\label{sec:discussion}

\subsection{Interpreting the Findings in Context}
Four methodological lessons emerge from this national isolated-CABG study.

First, \textbf{validation geometry matters far more than incremental model tuning.} Holding out calendar year 2017 mixes 99.5\% familiar hospitals with half-year support and documentation drift, and the 2.0-point AUC gap between naive 2017 and unseen 2017 hospitals dwarfs the separation among tuned GBDTs and frontier foundation models (\S\ref{sec:cohort_splits}). Reading published POAF AUCs alongside prediction window, surgical case mix, and predictor modality (\S\ref{sec:literature}) reconciles why pre-incision isolated-CABG models sit near 0.65--0.68 externally while later-window or multimodal tools exceed 0.85~\cite{pollock2018sts,fields2025poaf,ma2025noaf,khalpey2026digital}.

Second, \textbf{honest audits clarify where domain feature engineering helps.} Across 207 raw STS variables, adding 82 deterministic bedside composites acted as an inductive bias for full-registry multi-task neural networks ($+0.00102$, CI $>0$) and as a Stage-1 teacher distillation source ($+0.00057$, CI $>0$ into 207-raw EvoTarget-FT) rather than a standalone discrimination lift for 50k binary trees or foundation models, where AUC stayed flat and lower observed 15-bin ECE point estimates carried cluster intervals spanning zero (\S\ref{sec:agentic_fe}).

Third, \textbf{evolutionary continuous targets break the binary ceiling cleanly without late-stage collapse.} Extending target evolution across 1,000 generations (256,000 fits) proved that EWMA hardening plus 20\% confident-core gene-locking stabilizes smoothly rather than collapsing (0.67592 at Gen~1,000), peaking at Generation~20 where \textbf{EvoTarget-FT} reaches \textbf{0.676535} ($+0.001885$ $[+0.00026, +0.00358]$ vs.\ identical binary FT; 74.6\% of the 7$\times$ scaling gap) and lifting five distinct model paradigms with 74-hospital intervals strictly above zero. Our strictly nested 207-raw 2$\times$2 ablation further demonstrated that 65.2\% of this gain originates from pure zero-leakage target smoothing ($+0.001229$), whereas outer-fold global search contributes less than 1\,bp ($+0.000084$).

Fourth, \textbf{parallel symbolic/tree genetic programming complements clinical self-attention.} Re-engineering GPU-guided genetic programming (\textbf{JAIN Ultra GP v2}) erased six engineering bottlenecks to gain $+0.003984$ ($[+0.00196, +0.00601]$) over legacy JAIN (0.671666 $\to$ \textbf{0.675649}), beating every individual binary tabular ELO baseline with fresh CV-7 verified stability ($r=0.950$). Because sparse symbolic/linear$+$tree GP pipelines extract complementary log-odds signal to dense attention on $z^*$, their OOF-weighted \textbf{Dual-Paradigm Fusion} reaches \textbf{0.676677 ROC-AUC}, \textbf{0.377671 AP}, and \textbf{0.00724 ECE}---recovering \textbf{79.2\%} of the 7$\times$ scaling gap, strictly outperforming even a 3-way binary neural--GBDT--TabPFN blend ($+0.000992$ $[+0.00006, +0.00197]$), and exceeding 346k full-registry CatBoost on Average Precision, ECE, 2017 unseen-hospital transport, and rare clinical tails. Clinically, concentrating POAF risk to $\sim$44.9\% in the top decile (1.85$\times$ enrichment) offers a lean pre-incision enrichment tool; proving clinical benefit from model-guided prophylaxis requires prospective randomized trials~\cite{gaudino2021palacs,chatterjee2026sts}.

\subsection{Limitations and Path to Confirmation}
Five limitations bound the present claims under TRIPOD+AI and PROBAST+AI principles~\cite{collins2024tripodai,moons2025probastai}:
\begin{enumerate}
  \item \textbf{Retrospective selection exposure:} Although test labels never reached training, evolution, or GP search workers, the benchmark package evaluated 140+ model--target arms across the initial benchmark, 1,000-generation sweep, and GP v2 study before featuring EvoTarget-FT, JAIN Ultra GP v2, and Dual-Paradigm Fusion; hospital-cluster intervals condition on the selected runs~\cite{dwork2015reusable}. Confirmatory evidence requires freezing the pipeline prior to evaluating an unused cohort (e.g., STS ACSD v4.20+).
  \item \textbf{Residual teacher-feature distillation vs.\ pure raw-only power:} Our strictly nested 207-raw 2$\times$2 ablation resolved the non-nested outer-fold concern ($+0.000084$ global vs.\ nested gap) and isolated pure smoothing ($+0.001229$, 65.2\%) from 289$\to$207 teacher distillation ($+0.000572$, 30.3\%, CI $>0$). While the nested 207-raw point gain ($+0.001229$) is substantial, its standalone 74-hospital interval ($[-0.00033, +0.00285]$, $P>0=0.942$) still spans zero at $N=74$ centers, warranting larger multi-center confirmation.
  \item \textbf{Registry ascertainment and center heterogeneity:} STS \cli{COtAFib} captures treatment-requiring AF/flutter rather than all telemetry-detected episodes (33.4\% adjusted telemetry vs.\ 27.0\% STS across five centers~\cite{filardo2017telemetry,filardo2020telemetry}; 25\% vs.\ 15\% telemetry vs.\ daily ECG in isolated CABG~\cite{perezgrovas2023jaha}). Because center incidence spans 11.7\%--34.8\%, prospective deployment at a new hospital should pair frozen patient weights with lightweight local intercept recalibration~\cite{steyerberg2016prediction,riley2019external,vancalster2019calibration}.
  \item \textbf{Outcome-aware split design and data governance:} Both the 94.2/5.8 split search and 50k/10k class balancing used registry outcome counts, and STS v2.81 data are governed under institutional Data Use Agreement without patient-level re-operation linkage across hospitals.
  \item \textbf{Offline compute vs.\ deployment cost:} Offline Stage-1 Gen-20 discovery required 5,120 fits (and 256,000 fits to verify Gen-1,000 stability), whereas live inference averages only 20 compact transformer passes plus 30 lightweight symbolic/linear/tree GP passes using strictly 207 routine pre-incision variables.
\end{enumerate}

\subsection{Conclusion}
Reliable clinical prediction depends equally on evaluation hygiene, supervision quality, and paradigm diversity. Across 368,132 STS isolated-CABG procedures from 1,091 U.S.\ centers, we exposed why naive temporal splits misstate unseen-hospital transport, introduced leakage-audited hospital-disjoint benchmarks, clarified the scale-dependent value of 92 agent-engineered clinical features, and situated pre-incision STS discrimination within 42 published POAF studies. Against six world-leading tabular ELO baselines, 1,000-generation target evolution delivered hospital-cluster-confirmed gains across five paradigms and lifted EvoTarget-FT to 0.676535 (with 65.2\% verified under strictly nested 207-raw evolution), while re-engineered GPU-guided genetic programming (\textbf{JAIN Ultra GP v2}) jumped $+0.003984$ over legacy JAIN to 0.675649. Combining both paradigms in \textbf{Dual-Paradigm Fusion} achieved \textbf{0.676677 ROC-AUC}, \textbf{0.377671 AP}, and \textbf{0.00724 ECE}---recovering \textbf{79.2\%} of a 7$\times$ national data-scaling expansion on strictly 207 raw pre-incision predictors. Freezing this dual-paradigm pipeline for prospective confirmation on newer STS releases is the decisive next step toward reliable pre-incision POAF stratification.
